\documentclass[10pt,twocolumn]{article}
\usepackage[T1]{fontenc}
\usepackage[utf8]{inputenc}
\usepackage{lmodern}
\usepackage[margin=1.8cm,top=2.4cm,bottom=2cm,columnsep=0.6cm]{geometry}
\usepackage{fancyhdr}
\usepackage{titlesec}
\usepackage{booktabs}
\usepackage{amsmath,amssymb,amsfonts}
\usepackage{microtype}
\usepackage{xcolor}
\usepackage{mdframed}
\usepackage{parskip}
\usepackage{enumitem}
\usepackage{hyperref}

\definecolor{rulered}{RGB}{180,30,30}
\definecolor{boxgray}{RGB}{245,245,245}
\definecolor{keywordgray}{RGB}{100,100,100}

\pagestyle{fancy}
\fancyhf{}
\fancyhead[L]{\textsc{\small Claude Mag}}
\fancyhead[C]{\small\textit{Machine Learning Research Digest}}
\fancyhead[R]{\small 2026-07-20}
\fancyfoot[C]{\small\thepage}
\renewcommand{\headrulewidth}{0.8pt}

\titleformat{\section}{\normalsize\bfseries\color{rulered}}{}{0pt}{}%
  [\vspace{-4pt}\color{rulered}\rule{\columnwidth}{0.4pt}]
\titlespacing{\section}{0pt}{8pt}{4pt}

\hypersetup{colorlinks=true,urlcolor=rulered,linkcolor=black}

\begin{document}
\setlength{\parindent}{0pt}
\setlength{\parskip}{4pt}

{\small\color{keywordgray}\textbf{Keywords:}
  hallucination \textbullet{}
  parametric knowledge \textbullet{}
  evidence grounding \textbullet{}
  pretraining}\par\vspace{4pt}

{\LARGE\bfseries\raggedright
  Knowledgeless Language Models}\par\vspace{4pt}

{\large\itshape\raggedright
  Anonymizing Named Entities at Pretraining Time Suppresses
  Parametric Factual Recall While Improving Contextual QA,
  Fact Verification, and Hallucination Detection}\par\vspace{6pt}

{\small\textbf{By} Roi Cohen, Yvan Carr\'{e}, Nick Lechtenbörger,
  Hendrik Droste, Lucas Kerschke, Russa Biswas, Gerard de Melo,
  Jan Buys (Hasso Plattner Institute; University of Cape
  Town)}\par\vspace{2pt}

{\small\color{keywordgray}arXiv:2607.12831 \textbullet{}
  July 14, 2026}
\par\vspace{2pt}\noindent{\color{rulered}\rule{\columnwidth}{0.6pt}}\vspace{6pt}

Language models encode factual knowledge in their parameters during
pretraining, and this parametric knowledge is a double-edged sword.
It enables impressive zero-shot question answering but also produces
confident hallucinations when the model's memorized facts are
outdated, incomplete, or contradicted by the provided context.
KLLMs (Knowledgeless Language Models) address this at the source:
instead of correcting parametric recall through fine-tuning or prompting,
they prevent it from forming in the first place by anonymizing named
entities before pretraining begins.

\begin{mdframed}[backgroundcolor=boxgray,linecolor=rulered,linewidth=1pt,
    innerleftmargin=6pt,innerrightmargin=6pt,
    innertopmargin=6pt,innerbottommargin=6pt]
\textbf{\small AT A GLANCE}\par\vspace{4pt}
\begin{description}[leftmargin=0pt,labelsep=4pt,itemsep=2pt,parsep=0pt,
    font=\small\bfseries]
  \item[Problem:] LLMs memorize entity-linked facts during pretraining
    and hallucinate them confidently when context contradicts or
    omits them.
  \item[Why it matters:] Hallucination and parametric override are
    root causes of LLM unreliability in knowledge-intensive tasks;
    existing mitigations act at fine-tuning or inference time.
  \item[Method:] Named entities in the pretraining corpus are replaced
    with consistent-within-document, random-across-documents
    placeholder tokens, blocking entity-linked factual supervision.
  \item[Result:] +12.2\% relative on open-book NaturalQuestions;
    +8.5\% on FEVER; +9.3\% on HaluEval; near-zero perplexity cost.
  \item[Evidence:] 125M, 350M, and 1.3B model scales; five
    benchmarks spanning QA, fact verification, and hallucination
    detection.
\end{description}
\end{mdframed}

\section{The Big Picture}

When a language model is asked ``Who directed Oppenheimer?'' without
any context, it answers from its parameters. When it is asked the
same question with a context passage, it should answer from that
passage --- but it often does not, especially when the parametric
answer conflicts with the contextual one.

This \emph{parametric override} problem is distinct from the
out-of-distribution problem: it occurs even when the model has seen
the correct answer in context. The issue is that parametric knowledge
has been learned so strongly during pretraining that it dominates
contextual signals in downstream deployment.

KLLMs test a simple hypothesis: if named entities are removed from
pretraining data (replaced by anonymized placeholders), the model
cannot form entity-specific parametric associations in the first place,
making it structurally dependent on context for entity-linked facts.

\section{Why Existing Approaches Fall Short}

Retrieval-augmented generation (RAG) mitigates hallucination at
inference time but does not change the model's internal representation
or its tendency to override retrieved context with parametric beliefs.
Fine-tuning with contrastive or knowledge-editing objectives can
suppress specific facts but does not address the pretraining-level
memorization mechanism.

Prompting strategies (``base your answer only on the passage'')
improve context use in controlled settings but fail in adversarial
cases where the model's high-confidence parametric recall outweighs
the instruction. KLLMs target the fundamental source: the pretraining
signal that creates entity-linked parametric knowledge.

\section{Core Mechanism}

A named entity recognition (NER) model processes each document and
identifies spans corresponding to persons, organizations, locations,
dates, and named products. Each unique entity within a document is
replaced with a type-annotated placeholder: \texttt{[PERSON\_42]},
\texttt{[ORG\_7]}, etc.

Critically, the mapping from entities to placeholders resets between
documents. The same person gets placeholder \texttt{[PERSON\_42]} in
one document and \texttt{[PERSON\_19]} in another. This prevents
the model from learning cross-document entity associations ---
the mechanism by which parametric factual knowledge accumulates.

Within a document, coreference is preserved: if \texttt{[PERSON\_42]}
is the subject of the first sentence, it remains \texttt{[PERSON\_42]}
throughout, so the model still learns intra-document discourse structure.
Linguistic syntax, semantic roles, and factual relations within
each document are fully intact; only entity identity across documents
is suppressed.

\section{Technical Formulation}

Let $\mathcal{D} = \{d_1, d_2, \ldots, d_N\}$ be the pretraining
corpus. Let $E(d_i) = \{e_1, \ldots, e_k\}$ be the set of unique
entities in document $d_i$. The anonymized document is:
\[
\tilde{d}_i = \sigma_i(d_i)
\]
where $\sigma_i$ is a randomly sampled bijection from entities in
$d_i$ to type-specific placeholder tokens, drawn fresh for each
document. The KLLM is trained with the standard autoregressive
cross-entropy loss on $\{\tilde{d}_i\}$; no special objective is
needed.

The cross-document anonymization satisfies: for any two documents
$d_i \neq d_j$ and any entity $e$, with probability 1 the placeholder
for $e$ in $\tilde{d}_i$ differs from its placeholder in $\tilde{d}_j$.

\section{Experimental Findings}

\begin{center}
\small
\begin{tabular}{lcc}
\toprule
Benchmark & KLLM-1.3B & Baseline-1.3B \\
\midrule
NQ open-book EM$\uparrow$ & \textbf{43.1\%} & 38.4\% \\
NQ closed-book EM$\uparrow$ & 6.2\% & 18.7\% \\
FEVER accuracy$\uparrow$ & \textbf{83.7\%} & 79.2\% \\
HaluEval$\uparrow$ & \textbf{72.4\%} & 65.1\% \\
TruthfulQA$\uparrow$ & \textbf{48.3\%} & 41.7\% \\
Perplexity$\downarrow$ & 14.3 & 14.1 \\
\bottomrule
\end{tabular}
\end{center}

The closed-book NQ drop (18.7\%~$\to$~6.2\%) confirms that entity
anonymization successfully suppresses parametric factual recall.
The open-book NQ gain (+12.2\% relative) confirms that the model
compensates by better utilizing context. The perplexity difference
of 0.2 points (1.4\% relative) confirms that linguistic competence
is essentially unaffected.

\section{Ablations and Interpretation}

\textbf{Partial anonymization:} Anonymizing only person entities
(not organizations or locations) yields approximately half the HaluEval
improvement, suggesting that all entity types contribute to the
parametric recall problem.

\textbf{Fixed vs.\ random cross-document mapping:} Using a fixed
entity-to-placeholder mapping across all documents partially
re-introduces factual associations (the model can form cross-document
entity statistics by token co-occurrence). The random-per-document
approach is essential for full suppression.

\textbf{Downstream fine-tuning:} KLLMs fine-tuned on Wikipedia-based
QA recover most of the closed-book performance while retaining
the open-book advantage, suggesting that the intervention targets
pretraining memorization specifically rather than the model's capacity
to learn facts at all.

\textbf{Scale:} The open-book improvement grows from +6.1\% at 125M
to +12.2\% at 1.3B, suggesting that larger KLLMs are better at
leveraging context when parametric recall is suppressed.

\section*{Reference}

Roi Cohen et al.\
\textbf{Knowledgeless Language Models: Suppressing Parametric Recall
for Evidence-Grounded Language Modeling}.
arXiv:2607.12831, July 2026.
\url{https://arxiv.org/abs/2607.12831}

\end{document}
